\documentclass[10pt,a4paper]{amsart}
\usepackage[margin=19mm]{geometry}
\usepackage{amsmath,amssymb,mathtools}
\usepackage[scr=rsfs,cal=euler]{mathalfa}
\usepackage{xfrac}
\usepackage[unicode,hidelinks,bookmarks=false]{hyperref}
\newcommand{\ie}{i.e.\ }
\newcommand{\cf}{cf.\ }
\newcommand{\resp}{respectively\ }
\newcommand{\Subsection}{\subsection}
\providecommand{\nobreakdash}{\nobreak\hbox{-}}
\setlength{\emergencystretch}{2em}
\multlinegap0pt
\usepackage{bm}
\usepackage{bbm}
\usepackage{dsfont}
\usepackage{xspace}
\usepackage{cancel}
\usepackage{mathtools}
\usepackage{amsthm}
\makeatletter
\@ifundefined{proposition}{\newtheorem{proposition}{Proposition}}{}
\@ifundefined{theorem}{\newtheorem{theorem}{Theorem}}{}
\makeatother
\makeatletter
\newcommand{\R}{\mathbb{R}}
\newcommand{\bbt} {\mathbb T}
\newcommand{\caG}{\mathcal{G}}
\newcommand{\eps}{\varepsilon}
\expandafter\ifx\csname proofof\endcsname\relax
\newenvironment{proofof}[1]{{\flushleft {\it Proof of {#1}}}}
\fi
\expandafter\ifx\csname remark\endcsname\relax
\newenvironment{remark}{{\flushleft {\bf Remark.}}}{}
\fi
\expandafter\ifx\csname remarks\endcsname\relax
\newenvironment{remarks}{{\flushleft {\bf Remarks.}}}{}
\fi
\makeatother
\title[Decay of periodic entropy solutions to Euler-alignment syste]{Decay of periodic entropy solutions to Euler-alignment systems with non-constant kernel}
\author{Debora Amadori, Cleopatra Christoforou, Gianmarco Cipollone}
\date{Source: arXiv:2606.14446v1 (CC BY 4.0)}
\begin{document}
\maketitle

\noindent\emph{Source and licence.} Decay of periodic entropy solutions to Euler-alignment systems with non-constant kernel. Version arXiv:2606.14446v1 (CC BY 4.0), distributed under the Creative Commons Attribution 4.0 International licence, as recorded in the arXiv licence metadata for this exact version and shown on the abstract page https://arxiv.org/abs/2606.14446v1. Full rights notice: licenses/arxiv-2606.14446-CC-BY-4.0.txt.\par\smallskip

\noindent\textit{Editorial scope.} Counted unit: the Theorem whose source label is \texttt{th:energy-estimates} in the article (source0.tex lines 737--750, source label \texttt{th:energy-estimates}), delivered together with the article's own complete original proof of it (source0.tex lines 752--895, one \texttt{proof} environment of 144 lines). No mathematics is written, repaired or re-derived: statement and proof are the article's own text line for line. Required context retained uncounted: hyp-on-p (lines 211--213); hyp-on-K tilde (lines 384--386); eq:final-system (lines 404--409); eq:init-data tilde (lines 411--413); eq:bounds-rho tilde (lines 415--417); th:L2-decay (lines 448--454); def:H-function (lines 467--469); def:energy-function (lines 489--491); def:Z-function (lines 495--497); kappa-bound (lines 499--501); eq:property-G (lines 521--523); def:potential\_2 (lines 541--543); eq:partial-t.phi (lines 549--552); eq:partial-x-phi (lines 549--552); eq:bounds-phi (lines 562--564); def:c1c2 (lines 610--614); eq:entropy-inequality (lines 618--620); bound-on-sigma\_1 (lines 648--650); prop:3.3 (lines 696--709); bound-on-sigma\_2 (lines 698--700); eq:constants-c (lines 706--708). Every definition, display and label of those lines is retained unchanged. Omitted from this package and not redistributed: 1--210, 214--383, 387--403, 410--410, 414--414, 418--447, 455--466, 470--488, 492--494, 498--498, 502--520, 524--540, 544--548, 553--561, 565--609, 615--617, 621--647, 651--695, 701--705, 709--736, 751--751, 896--1006. Nothing inside a retained excerpt is omitted or altered: every retained body is delivered exactly as the article's lines are written, so no omission receipt of any kind accompanies this package. Citations: the retained excerpts carry no citation command at all, so no bibliography entry of the article is transcribed and no reference is redistributed. Reference adaptation: none was needed and none was made; every reference inside the delivered unit resolves inside it, so no unresolved reference or citation is produced. Printed slips kept literal: no printed word, symbol, hypothesis or number of the counted statement or of its proof was altered. Layout and class: the article's own class is \texttt{amsart} with packages including hyperref, amssymb, mathrsfs, graphicx, subfigure, amsfonts, amsmath, enumerate, amsthm, colortbl, xcolor, floatrow, soul, comment; the wrapper is the lane's pinned \texttt{amsart} editorial wrapper at 10pt on A4, which restates the theorem-like environments the retained text needs and adds the article's own macro definitions that the retained text needs (\texttt{\textbackslash R}, \texttt{\textbackslash bbt}, \texttt{\textbackslash caG}, \texttt{\textbackslash eps}), with extra wrapper packages (bm, bbm, dsfont, xspace, cancel, mathtools). The delivered numbering is the unit's own. Assets: no retained span contains a figure environment, a graphics-inclusion command, a PDF-page inclusion, a table or a TikZ picture; no image file is copied into this package, the package declares no asset dependency and no image file is redistributed with it. Licence: version arXiv:2606.14446v1 of this article is distributed under the Creative Commons Attribution 4.0 International licence, as recorded in the arXiv licence metadata for this exact version and shown on the abstract page https://arxiv.org/abs/2606.14446v1. A full rights notice is delivered at licenses/arxiv-2606.14446-CC-BY-4.0.txt. Characters: every retained excerpt is pure ASCII, so the delivered engine, XeLaTeX with its default Latin Modern text font, reports no missing character.
\begin{equation}\label{hyp-on-p}
  p\in C^1([0,+\infty);[0,+\infty))\,,\qquad  p'>0
\end{equation}

\begin{equation}\label{hyp-on-K tilde}
\tilde K: \bbt\to \R\,,\qquad \tilde K \in L^1(\bbt)\,,\qquad \tilde K \mbox{ even}\,,\qquad \tilde K(x)\ge 0\,.
\end{equation}

\begin{equation}\label{eq:final-system}
    \begin{cases}
        \partial_t\tilde\rho+\partial_x \tilde m=0,\\
        \partial_t \tilde m+ \partial_x \left(\frac{\tilde m^2}{\tilde\rho+\bar\rho}+p(\tilde\rho+\bar\rho)\right) = -2\bar K\bar\rho \, \tilde m(x) + \tilde \caG[\tilde \rho+\bar\rho, \tilde m]
    \end{cases}
\end{equation}

\begin{equation}\label{eq:init-data tilde}
(\tilde\rho,\tilde m)(x,0)=\left(\tilde\rho_0(x), \tilde m_0(x)\right)\,\qquad x\in\bbt\;,
\end{equation}

\begin{equation}\label{eq:bounds-rho tilde}
  - \bar\rho<\bar\rho_{inf}-\bar\rho \le \tilde\rho(x,t)\le \bar\rho_{sup} -\bar\rho
\end{equation}

\begin{theorem}\label{th:L2-decay}  %Suppose that 
Assume \eqref{hyp-on-p}, \eqref{hyp-on-K tilde} and that $\bar K>0$, $\bar\rho>0$. Furthermore, assume there exists an entropy weak solution $(\tilde\rho,\tilde m)(\cdot,t)$ to problem \eqref{eq:final-system}--\eqref{eq:init-data tilde} for $t\in [0,T^*)$ that satisfies~\eqref{eq:bounds-rho tilde} on $\Omega$ for some $\bar\rho_{inf}<\bar\rho<\bar\rho_{sup}$. Then there exist positive constants $\mu$ and $c$ such that
    \begin{equation}\label{prop:three}
         \| (\tilde \rho,\tilde m)(\cdot,t) \|_{L^2(\bbt)} \leq \mu\, e^{-ct}\, \| (\tilde \rho_0,\tilde m_0)\|_{ L^2(\bbt)},\qquad \forall \, t\in [0,T^*)
    \end{equation}
    with $\mu$ and $c$ independent of time. 
\end{theorem}

\begin{align}
    \mathcal{H}(t) & := \frac{1}{2}\int_\bbt\int_\bbt \tilde K(x-x')\left[\tilde m(x,t)\sqrt{\frac{\tilde\rho(x',t)+\bar\rho}{\tilde\rho(x,t)+\bar\rho}}-\tilde m(x',t)\sqrt{\frac{\tilde\rho(x,t)+\bar\rho}{\tilde\rho(x',t)+\bar\rho}} \,\right]^2dx\,dx'\,.\label{def:H-function}
    \end{align}

    \begin{equation} \label{def:energy-function}
    \mathcal{E}(t) := \int_\bbt \left[ \sigma\tilde\eta + \Phi^2 - \frac{1}{\bar K\bar\rho}\Phi \tilde m \right]\,dx 
    \end{equation}

    \begin{align}
        \mathcal{Z}(t) & := \left(\sigma\kappa-\frac{1}{\bar K}\right)\int_\bbt \frac{\tilde m^2}{\tilde\rho+ \bar\rho}\,dx + \beta \int_\bbt \tilde\rho^2\,dx,\label{def:Z-function}
\end{align}

\begin{equation}\label{kappa-bound}
 0<   \kappa < 2\bar K\bar\rho\,,\qquad 0<\beta<  \frac{C}{\bar K\bar\rho}\,,
\end{equation}

 \begin{equation}\label{eq:property-G}
        \mathcal{H}(t)=-\int_\bbt \frac{\tilde m(x)}{\tilde \rho(x)+\bar\rho}\tilde\caG[\tilde\rho+\bar\rho, \tilde m]\,dx \ge 0\;,
\end{equation}

\begin{equation}\label{def:potential_2}
    \Phi(x,t):= \psi(x,t)- <\psi(\cdot,t)>\,,      
\end{equation}

\begin{align}
    \partial_x\Phi(x,t) & = \tilde\rho(x,t),\label{eq:partial-x-phi}\\
    \partial_t\Phi(x,t) & = - \tilde m(x,t).\label{eq:partial-t.phi}
\end{align}

\begin{equation}\label{eq:bounds-phi}
    \int_{-1}^1 \Phi^2(x,t)\,dx \leq 2\int_{-1}^1\tilde \rho^2(x,t)\,dx\;,\qquad\forall t\in[0,T^*)\,.
\end{equation}

\begin{equation}\label{def:c1c2}
c_1 = \min\left\{\min_{[\bar\rho_{inf},\bar\rho_{sup}]}\frac{p'(\rho)}{\rho},\,\frac{1}{2 \bar\rho_{sup} }\right\}
\,,\qquad
c_2 = \max\left\{\max_{[\bar\rho_{inf},\bar\rho_{sup}]}\frac{p'(\rho)}{\rho},\,\frac{1}{2 \bar\rho_{inf} }\right\}.
\end{equation} 

\begin{equation}\label{eq:entropy-inequality}
    \partial_t\tilde\eta(\tilde\rho,\tilde m)+\partial_x\tilde q(\tilde\rho,\tilde m) \leq \frac{\tilde m}{\tilde\rho+\bar\rho}\left( -2\bar K\bar\rho\, \tilde m + \tilde \caG[\tilde\rho+\bar\rho,\tilde m] \right).
\end{equation}

\begin{equation}\label{bound-on-sigma_1}
\sigma> \sigma_1:=\frac{1}{c_1\bar K\bar\rho}    \;.
\end{equation} 

\begin{proposition}\label{prop:3.3}
If \eqref{def:c1c2}, \eqref{kappa-bound},~\eqref{bound-on-sigma_1} hold and 
\begin{equation}\label{bound-on-sigma_2}
\sigma>\sigma_2:=\frac{1}{\kappa\bar K}\,,
\end{equation}
then 
\begin{equation}\label{Z-E}
    \mathcal Z(t)\geq 2c \,\mathcal E(t)\,,\qquad t\in[0,T^*)
\end{equation}
where
\begin{equation}
c :=\min\left\{\frac{1}{\bar\rho_{sup}}\frac{\bar\rho(\sigma\kappa \bar K -1)}{1+2\bar K\bar\rho c_2\sigma},\frac{1}{2}\frac{\beta\bar K\bar\rho}{\bar K\bar\rho c_2\sigma+2\bar K\bar\rho+1} \right\} >0\;.\label{eq:constants-c}
\end{equation}
\end{proposition}

\begin{theorem}\label{th:energy-estimates}
    Under the assumptions of Theorem \ref{th:L2-decay}, %there exists
    %positive constants $c,\,c_1,\,c_2,\alpha$ and 
if $\sigma$ is chosen sufficiently large and \eqref{kappa-bound} hold,
   then the functionals $\mathcal H$, $\mathcal E$ and $\mathcal Z$ given at~\eqref{def:H-function}, \eqref{def:energy-function} and \eqref{def:Z-function} respectively, satisfy
    \begin{equation}\label{eq:inequality-EHZ1}
        \mathcal E(t_2)-\mathcal E(t_1) + \int_{t_1}^{t_2} \left[\sigma \mathcal H(t) +  \mathcal Z(t) \right]\,dt \leq 0
    \end{equation}
and  
    \begin{equation}\label{eq:gronw-energy}
        \mathcal{E}(t) + \sigma\int_0^t \mathcal H(s)e^{-2c(t-s)}\,ds \leq \mathcal E(0)e^{-2ct}
    \end{equation}
for any $t,\,t_1,\,t_2\in[0, T^*)$ with $t_2>t_1$, where $c$ is given in \eqref{eq:constants-c}. 
\end{theorem}

\begin{proof} 
To begin with, we multiply the second equation of system~\eqref{eq:final-system} by the potential $\Phi(t)$, 
then add the term 
$\tilde m\partial_t\Phi+\left( \frac{\tilde m^2}{\tilde\rho+\bar\rho}+p(\tilde\rho+\bar\rho) \right)\partial_x\Phi$ 
to get
\begin{align}
    &\partial_t(\Phi \tilde m)+\partial_x \left(\Phi\left(\frac{\tilde m^2}{\tilde \rho+\bar\rho}+p(\tilde\rho+\bar\rho)\right)\right) \nonumber\\
    &\qquad =\Phi(-2\bar K \bar\rho \,\tilde m + \tilde \caG[\tilde\rho+\bar\rho,\tilde m])-\tilde m^2+\tilde\rho\left( \frac{\tilde m^2}{\tilde\rho+\bar\rho}+p(\tilde\rho+\bar\rho) \right)\nonumber\\
    &\qquad = -2\Phi\bar K \bar\rho \tilde m + \Phi\tilde \caG[\tilde\rho+\bar\rho,\tilde m] -\tilde m^2 \frac{\bar \rho}{\tilde\rho+\bar\rho}+\tilde\rho \,p(\tilde\rho+\bar\rho)\label{eq:Phi-m}\,,
\end{align}
using \eqref{eq:partial-x-phi}--\eqref{eq:partial-t.phi}. Furthermore, multiplying \eqref{eq:partial-t.phi} by $2\Phi$, we obtain 
\begin{equation}\label{eq:Phi^2}
\partial_t\Phi^2 = -2\tilde m\,\Phi.
\end{equation}
Now, we write inequality \eqref{eq:entropy-inequality} and equations~\eqref{eq:Phi-m} and \eqref{eq:Phi^2} in distributional sense. 
Indeed, for any nonnegative periodic test function $\varphi\in C^1(\bbt\times (0,T^*))$ with support contained in $\bbt\times [t_1,t_2]$ 
for some $0<t_1<t_2<T^*$, we combine these three integral relations to reach the inequality 
\begin{align}
    \int_0^{T^*}\int_\bbt & \left(\sigma\tilde\eta +\Phi^2 -\frac{1}{\bar K\bar\rho}\Phi \,\tilde m\right)\partial_t\varphi\,dx\,dt  \nonumber\\ 
    & \geq \int_0^{T^*}\int_\bbt -\sigma\tilde q\partial_x\varphi +\sigma\frac{\tilde m}{\tilde\rho+\bar\rho}\left( 2\bar K\bar\rho \,\tilde m -\tilde \caG[\tilde\rho+\bar\rho,\tilde m] \right)\varphi\,dx\,dt\nonumber\\
    & \quad+\int_0^{T^*}\int_\bbt \frac{1}{\bar K\bar\rho}\left( \Phi\left(\frac{\tilde m}{\tilde\rho+\bar\rho}+p(\tilde\rho+\bar\rho)\right) \right)\partial_x \varphi \,dx\,dt\nonumber\\
    &\quad+\int_0^{T^*}\int_\bbt -\frac{1}{\bar K\bar\rho}\left( 2\Phi\bar K\bar\rho \,\tilde m - \Phi\,\tilde \caG[\tilde\rho+\bar\rho, \tilde m] + \tilde m^2\frac{\bar\rho}{\tilde\rho+\bar\rho} -\tilde\rho \,p(\tilde\rho+\bar\rho)\right)\varphi\,dx\,dt\nonumber\\
    &\quad +\int_0^{T^*}\int_\bbt 2\tilde m\,\Phi\,\varphi\,dx\,dt\,.\label{eq:ineq-energy-1}
\end{align}
Notice that the right hand side of the inequality above is equal to
\begin{align}
    \int_0^{T^*}\int_\bbt & \left[ \frac{1}{\bar K\bar\rho}\left( \Phi\frac{\tilde m}{\tilde\rho+\bar\rho}+\Phi  \,p(\tilde\rho+\bar\rho) \right) - \sigma \tilde q\right]\partial_x\varphi\,dx\,dt+\left(2\sigma\bar K\bar\rho-\frac{1}{\bar K}\right)\int_0^{T^*}\int_\bbt \frac{\tilde m^2}{\tilde\rho+\bar\rho}\varphi\,dx\,dt\nonumber\\
    &\quad+ \int_0^{T^*}\int_\bbt\left[ \frac{1}{\bar K\bar\rho} \left( \Phi\,\tilde \caG[\tilde\rho+\bar\rho,\tilde m]+\tilde\rho\, p(\tilde\rho+\bar\rho) \right) - \sigma\frac{\tilde m}{\tilde\rho+\bar\rho}\tilde \caG[\tilde\rho+\bar\rho,\tilde m]\right]\varphi\,dx\,dt\,.\label{eq:distrib-ineq-energy}
\end{align}
Given $0<t_1<t_2<T^*$, let $\eps>0$ be small such that $t_1-\eps\geq0$.
We now select the test function to be  $\varphi(x,t)=\chi(x)h_\eps(t)\in C^1(\bbt\times (0,T^*))$ with
\begin{align*}
    \chi(x) & :=1,\qquad x\in\bbt\\
    h_\eps(t) & := 
    \begin{cases}
        0, & 0\leq t<t_1-\eps,\\
        \eps^{-1}(t-t_1)+1,\quad  &  t_1-\eps \leq t < t_1,\\
        1, & t_1 \leq t \leq t_2,\\
        \eps^{-1}(t_2-t) + 1, & t_2 < t \leq t_2+\eps,\\
        0, & t> t_2+\eps,
    \end{cases}\qquad t\in (0,T^*)
\end{align*}
and evaluate~\eqref{eq:distrib-ineq-energy} for the test function $\phi(x,t)=\chi(x)h_\eps(t)$ to get
\begin{align*}
    & \left( 2\sigma\bar K\bar\rho-\frac{1}{\bar K}\right)\int_{t_1-\eps}^{t_2+\eps}\int_\bbt \frac{\tilde m^2}{\tilde\rho+\bar\rho}h_\eps(t)\,dx\,dt\\
    & \qquad+ \int_{t_1-\eps}^{t_2+\eps}\int_\bbt \left[ \frac{1}{\bar K\bar\rho}\left( \Phi\, \tilde \caG[\tilde\rho+\bar\rho, \tilde m] + \tilde\rho\, p(\tilde\rho+\bar\rho) \right) - \sigma\frac{\tilde m}{\tilde\rho+\bar\rho}\tilde \caG [\tilde\rho+\bar\rho, \tilde m] \right]h_\eps(t)\,dx\,dt.
\end{align*}
Using the definition~\eqref{def:energy-function} of $\mathcal{E}(t)$, inequality~\eqref{eq:ineq-energy-1}  and the above computations, we now let $\eps\to 0+$ to obtain
\begin{align}
    \mathcal{E}(t_2) & -\mathcal{E}(t_1) \leq  -\left( 2\sigma\bar K\bar\rho-\frac{1}{\bar K}\right)\int_{t_1}^{t_2}\int_\bbt \frac{\tilde m^2}{\tilde\rho+\bar\rho}\,dx\,dt\nonumber\\
    & \quad- \int_{t_1}^{t_2}\int_\bbt \left[ \frac{1}{\bar K\bar\rho}\left( \Phi \,\tilde \caG[\tilde\rho+\bar\rho, \tilde m] +\tilde \rho \,p(\tilde\rho+\bar\rho) \right) - \sigma\frac{\tilde m}{\tilde\rho+\bar\rho}\tilde \caG [\tilde\rho+\bar\rho,\tilde m] \right]\,dx\,dt.\nonumber
\end{align}
The last inequality can be rewritten as
\begin{align}\label{eq:en H ineq}
    \mathcal E(t_2)-\mathcal E(t_1) + \sigma\int_{t_1}^{t_2} \mathcal H(t)\,dt \leq -\left( 2\sigma\bar K\bar\rho-\frac{1}{\bar K}\right)\int_{t_1}^{t_2}\int_\bbt \frac{\tilde m^2}{\tilde\rho+\bar\rho}\,dx\,dt\nonumber\\
    - \int_{t_1}^{t_2}\int_\bbt \left[ \frac{1}{\bar K\bar\rho}\left( \Phi \tilde \caG[\tilde\rho+\bar\rho, \tilde m] +\tilde \rho \,p(\tilde\rho+\bar\rho) \right)\right]\,dx\,dt\;,
\end{align}
using~\eqref{eq:property-G} for $\rho=\tilde\rho+\bar\rho$ and $m=\tilde m$.

Recalling \eqref{hyp-on-p}, we can use Taylor expansion with Lagrange remainder
$$p(\tilde\rho+\bar\rho) = p(\bar\rho)+ p'(\xi)\tilde\rho    
$$
for some $\xi\in (\bar\rho_{inf},\bar\rho_{sup})$. Hence 
$$\tilde\rho p(\tilde\rho+\bar\rho) =\tilde\rho p(\bar\rho)+ p'(\xi)\tilde\rho^2 \ge \tilde\rho p(\bar\rho)+ C\tilde\rho^2\;,
$$
with 
$$C:= \inf_{r\in (\bar\rho_{inf},\bar\rho_{sup})} p'(r) >0\;.$$
Since $\int_{\bbt}\tilde \rho \,dx=0$, then
$$\int_\bbt \tilde\rho p(\tilde\rho+\bar\rho)\,dx \ge C \int_\bbt\tilde \rho^2 \,dx,
$$
and in view of the above analysis, inequality~\eqref{eq:en H ineq} takes the form
\begin{align}\label{eq:E-H}
    \mathcal E(t_2)-\mathcal E(t_1) + \sigma\int_{t_1}^{t_2} \mathcal H(t)\,dt \leq -\left( 2\sigma\bar K\bar\rho-\frac{1}{\bar K}\right)\int_{t_1}^{t_2}\int_\bbt \frac{\tilde m^2}{\tilde\rho+\bar\rho}\,dx\,dt\nonumber\\
    - \int_{t_1}^{t_2}\int_\bbt \left[ \frac{1}{\bar K\bar\rho}\left( \Phi \,\tilde \caG[\tilde\rho+\bar\rho,\tilde m] + C \tilde\rho^2 \right)\right]\,dx\,dt.
\end{align}
Next, by means of the definition of $\mathcal{Z}$ in \eqref{def:Z-function} we find
\begin{align}
    & \mathcal{E}(t_2)-\mathcal{E}(t_1) + \int_{t_1}^{t_2}\left[\sigma \mathcal{H}(t)+ \mathcal{Z}(t)\right]\,dt \nonumber\\
   & \qquad \le 
       - \sigma(2\bar K\bar\rho-\kappa)\int_{t_1}^{t_2}\int_\bbt\frac{\tilde m^2}{\tilde\rho+\bar\rho}\,dx  + \left(\beta -\frac{C}{\bar K\bar\rho}\right)\int_{t_1}^{t_2}\int_\bbt \tilde\rho^2\,dx \nonumber \\
        &\qquad \quad - \frac{1}{\bar K\bar\rho}\int_{t_1}^{t_2}\int_\bbt\Phi\, \tilde \caG(\tilde\rho+\bar\rho,   \tilde m)\,dx.\label{E+sH+Z}
\end{align}
Because of \eqref{kappa-bound}, we know that the first two terms are negative. 
Hence, it remains to control 
the last term $-\frac{1}{\bar K\bar\rho}\int_\bbt\Phi(x)\,\tilde \caG[\tilde\rho+\bar\rho,\tilde m](x)\,dx$. 
This one can be split into two integrals as follows:
\begin{align*}
     - & \frac{1}{\bar K\bar\rho} \int_\bbt\Phi(x)\tilde \caG[\tilde\rho+\bar\rho,\tilde m](x)\,dx \\
    & = -\frac{1}{\bar K\bar\rho}\int_\bbt\int_\bbt \tilde K (x-x')\Phi(x)[\tilde\rho(x) \tilde m(x')-\tilde\rho(x')\tilde m(x)]\,dx\,dx'\\
    &\qquad-\frac{1}{\bar K}\int_\bbt\int_\bbt \tilde K (x-x')\Phi(x)[\tilde m(x') - \tilde m(x)]\,dx\,dx'=:J_1+J_2
\end{align*}
Noting that, by the symmetry of the integrand, it holds
\begin{equation*}
    \frac{1}{\bar K\bar\rho}\int_\bbt\int_\bbt \tilde K(x-x')[\tilde \rho(x)\tilde m(x')-\tilde\rho(x')\tilde m(x)]\,dxdx'=0, 
\end{equation*}
then from~\eqref{def:potential_2}, we can reduce term $J_1$ to
\begin{align*}
J_1=- & \frac{1}{\bar K\bar\rho}\int_\bbt\int_\bbt \tilde K (x-x')\psi(x)[\tilde\rho(x) \tilde m(x')-\tilde\rho(x') \tilde m(x)]\,dx\,dx'\;.
\end{align*}
By estimating, we get 
\begin{align*}
 J_1=   -\frac{1}{\bar K\bar\rho} & \int_\bbt\int_\bbt \tilde K (x-x')\psi(x)[\tilde \rho(x) \tilde m(x')-\tilde \rho(x')\tilde m(x)]\,dx\,dx'\\
    \leq & \frac{1}{\bar K\bar\rho}\int_\bbt\int_\bbt |\psi(x)|\tilde K(x-x')|\tilde\rho(x) \tilde m(x') - \tilde\rho(x') \tilde m(x)|\,dxdx'\\
    \leq & \frac{1}{\bar K\bar\rho}\int_\bbt\int_\bbt (\bar\rho_{sup}-\bar\rho)\tilde  K(x-x')(|\tilde \rho(x) \tilde m(x')|+|\tilde\rho(x') \tilde m(x)|)\,dxdx'\\
    = & \frac{2(\bar\rho_{sup}-\bar\rho)}{\bar K\bar\rho}\int_\bbt\int_\bbt \tilde K(x-x')|\tilde\rho(x) \tilde m(x')|\,dxdx'\\
    \leq & \frac{(\bar\rho_{sup}-\bar\rho)}{\bar K\bar\rho}\int_\bbt\int_\bbt \tilde K(x-x')[\frac{1}{\alpha^2}\tilde\rho^2(x)+\alpha^2 \tilde m^2(x')]\,dxdx'\\
    = & \frac{\|\tilde K\|_{L^1}(\bar\rho_{sup}-\bar\rho)}{\bar K\bar\rho}\int_\bbt[\frac{1}{\alpha^2}\tilde\rho^2(x)+\alpha^2 \tilde m^2(x)]\,dx,
\end{align*}
for a parameter $\alpha>0$ to be determined. About $J_2$, we proceed in a similar fashion and also use~\eqref{eq:bounds-phi} to obtain
\begin{align*}
    J_2 =~&-\frac{1}{\bar K}  \int_\bbt\int_\bbt \tilde K (x-x')\Phi(x)[\tilde m(x') - \tilde m(x)]\,dx\,dx'\\ 
    =~& -\frac{1}{\bar K}\int_\bbt\int_\bbt \tilde K(x-x')\Phi(x)\tilde m(x')\,dxdx'
    %\\
   %& \qquad \qquad 
   + \frac{1}{\bar K}\int_\bbt\int_\bbt \tilde K(x-x')\Phi(x)\tilde m(x)\,dxdx'\\
     \leq ~& \frac{1}{2\bar K}\int_\bbt\int_\bbt \tilde K(x-x')[\frac{1}{\alpha^2}\Phi^2(x)+\alpha^2\tilde m^2(x')]\,dxdx'\\
    &\qquad \qquad + \frac{1}{2\bar K}\int_\bbt\int_\bbt \tilde K(x-x')[\frac{1}{\alpha^2}\Phi^2(x)+\alpha^2\tilde m^2(x)]\,dxdx'\\
   =~ &\frac{1}{\bar K\alpha^2}\int_\bbt\int_\bbt \tilde K(x-x')\Phi^2(x)\,dxdx' + \frac{\alpha^2}{\bar K}\int_\bbt\int_\bbt \tilde K(x-x') \tilde m^2(x)\,dxdx'\\
    =  ~&  \frac{\| \tilde K\|_{L^1}}{\bar K\alpha^2}\int_\bbt\Phi^2(x)\,dx + \frac{\alpha^2\| \tilde K\|_{L^1}}{\bar K}\int_\bbt  \tilde m^2(x)\,dx\\
  \leq ~&  \frac{2\| \tilde K\|_{L^1}}{\alpha^2\bar K}\int_\bbt\tilde\rho^2(x)\,dx + \frac{\alpha^2\| \tilde K\|_{L^1}}{\bar K}\int_\bbt  \tilde m^2(x)\,dx.
\end{align*}
Combining the above with estimate~\eqref{E+sH+Z}, we arrive at
\begin{align}
    \mathcal{E}(t_2)  &-\mathcal{E}(t_1)  + \int_{t_1}^{t_2}\left[\sigma \mathcal{H}(t) + \mathcal{Z}(t)\right]\,dt \nonumber\\
    & \leq \left[- \frac{\sigma(2\bar K\bar\rho-\kappa)}{\bar\rho_{sup}}+\frac{\|\tilde K\|_{L^1}}{\bar K}\alpha^2\left(  1+\frac{\bar\rho_{sup}-\bar\rho}{\bar\rho}\right)\right]\int_{t_1}^{t_2}\int_\bbt \tilde m^2\,dx \nonumber\\
    & \qquad + \left[\frac{\|\tilde K\|_{L^1}}{\bar K\alpha^2}\left( 2+\frac{\bar\rho_{sup}-\bar\rho}{\bar\rho} \right) + \beta - \frac{C}{\bar K\bar\rho}\right]\int_{t_1}^{t_2}\int_\bbt\tilde\rho^2\,dx\,. \label{eq:E+sigma_H+Z}
\end{align}
Having $\kappa$ and $\beta$ satisfying~\eqref{kappa-bound}, we choose $\alpha$ and $\sigma$ 
so that the right hand side of inequality~\eqref{eq:E+sigma_H+Z} to be non-positive. 
So, we first select $\alpha>0$ to fulfill
\begin{equation}\label{cond-on-alpha}
\alpha^2\geq \frac{\| \tilde K \|_{L^1(\bbt)}(\bar\rho+\bar\rho_{sup})}{C-\beta\bar K\bar\rho}
\end{equation}
and then $\sigma>\sigma_3$, with
\begin{equation}\label{bound-on-sigma_3}
\sigma_3:=\alpha^2 \frac{\| \tilde K \|_{L^1(\bbt)}}{\bar\rho\bar K}\frac{\bar\rho_{sup}^2}{2\bar K\bar\rho-\kappa}\;.
\end{equation}
Thus,~\eqref{eq:inequality-EHZ1} holds true. Moreover, thanks to Proposition \ref{prop:3.3}, it holds true~\eqref{eq:gronw-energy} by Gronwall Lemma. 
In summary, if $\sigma>\max\{\sigma_1,\sigma_2,\sigma_3\}$ given at ~\eqref{bound-on-sigma_1}, \eqref{bound-on-sigma_2} and \eqref{bound-on-sigma_3}, that is
\begin{equation}
\sigma  > \max\left\{ \alpha^2 \frac{\| \tilde K \|_{L^1(\bbt)}}{\bar\rho\bar K}\frac{\bar\rho_{sup}^2}{2\bar K\bar\rho-\kappa}, \frac 1 {\kappa \bar K}, \frac{1}{c_1\bar K\bar\rho}\right\},\label{eq:constants}    
\end{equation}
then both~\eqref{eq:inequality-EHZ1} and ~\eqref{eq:gronw-energy} follow. The proof of Theorem~\ref{th:energy-estimates} is complete.
\end{proof}

\end{document}
